

\documentclass{article}

% -----------------------------
% Packages
% -----------------------------
\usepackage{fancyhdr}
\usepackage{extramarks}
\usepackage{amsmath}
\usepackage{amsthm}
\usepackage{amsfonts}
\usepackage{tikz}
\usepackage[plain]{algorithm}
\usepackage{algpseudocode}
\usepackage{amssymb}
\usepackage{enumitem}

\usetikzlibrary{automata,positioning}

% -----------------------------
% Basic Document Settings
% -----------------------------
\topmargin=-0.45in
\evensidemargin=0in
\oddsidemargin=0in
\textwidth=6.5in
\textheight=9.0in
\headsep=0.25in

\linespread{1.1}

\pagestyle{fancy}
\lhead{\hmwkAuthorName}
\chead{\hmwkClass: \hmwkTitle}
\rhead{}
\lfoot{\lastxmark}
\cfoot{\thepage}

\renewcommand\headrulewidth{0.4pt}
\renewcommand\footrulewidth{0.4pt}

\setlength\parindent{0pt}

% -----------------------------
% Problem Section Helpers
% -----------------------------
\newcommand{\enterProblemHeader}[1]{
  \nobreak\extramarks{}{Problem #1 continued on next page\ldots}\nobreak{}
  \nobreak\extramarks{Problem #1 (continued)}{Problem #1 continued on next page\ldots}\nobreak{}
}

\newcommand{\exitProblemHeader}[1]{
  \nobreak\extramarks{Problem #1 (continued)}{Problem #1 continued on next page\ldots}\nobreak{}
  \nobreak\extramarks{Problem #1}{}\nobreak{}
}

\setcounter{secnumdepth}{0}
\newcounter{partCounter}
\newcounter{homeworkProblemCounter}
\setcounter{homeworkProblemCounter}{1}
\nobreak\extramarks{Problem \arabic{homeworkProblemCounter}}{}\nobreak{}

% -----------------------------
% Theorem env (optional)
% -----------------------------
\newtheorem*{theorem}{Theorem}

% -----------------------------
% Homework Problem Environment
%   Usage:
%   \begin{homeworkProblem}            % auto-number
%     % Your answer here
%   \end{homeworkProblem}
%
%   \begin{homeworkProblem}[1.2.3]     % label problem "1.2.3" (e.g. textbook numbering)
%     % Your answer here
%   \end{homeworkProblem}
% -----------------------------
\newenvironment{homeworkProblem}[1][]{
  \def\hwProblemLabel{#1}
  \ifx\hwProblemLabel\empty
    \edef\hwProblemLabel{\arabic{homeworkProblemCounter}}
    \stepcounter{homeworkProblemCounter}
  \fi
  \section{Problem \hwProblemLabel}
  \setcounter{partCounter}{1}
  \enterProblemHeader{\hwProblemLabel}
}{
  \exitProblemHeader{\hwProblemLabel}
}

% -----------------------------
% Homework Details (EDIT THESE)
% -----------------------------
\newcommand{\hmwkTitle}{Homework \#9}
\newcommand{\hmwkClass}{Math 3320}
\newcommand{\hmwkAuthorName}{\textbf{Joseph Quinn}}

% -----------------------------
% Title
% -----------------------------
\title{
  \vspace{2in}
  \textmd{\textbf{\hmwkClass:\ \hmwkTitle}}\\
  \vspace{3in}
}
\author{\hmwkAuthorName}
\date{}

% -----------------------------
% Convenience Macros (optional)
% -----------------------------
\renewcommand{\part}[1]{\textbf{\large Part \Alph{partCounter}}\stepcounter{partCounter}\\}

% Algorithms
\newcommand{\alg}[1]{\textsc{\bfseries \footnotesize #1}}

% Calculus
\newcommand{\deriv}[1]{\frac{\mathrm{d}}{\mathrm{d}x} (#1)}
\newcommand{\pderiv}[2]{\frac{\partial}{\partial #1} (#2)}
\newcommand{\dx}{\mathrm{d}x}

% Statistics
\newcommand{\E}{\mathrm{E}}
\newcommand{\Var}{\mathrm{Var}}
\newcommand{\Cov}{\mathrm{Cov}}
\newcommand{\Bias}{\mathrm{Bias}}

% Proof step (for aligned derivations)
\newcommand{\step}[2]{& #1 & & \text{#2} \\}

% Solution header (optional)
\newcommand{\solution}{\textbf{\large Solution}}

% -----------------------------
% Document
% -----------------------------
\begin{document}

\begin{homeworkProblem}[A17]
	We are given
	\[
		x_is_i+x_i\Delta s_i+s_i\Delta x_i=0,
	\]
	where $x_i>0$ and $s_i>0$.

	\bigskip

	First lets suppose $\Delta x_i>0$. we want to prove $s_i'<0$.

	Starting with the given equation,
	\[
		\begin{aligned}
			x_is_i+x_i\Delta s_i+s_i\Delta x_i & =0  \\
			x_i(s_i+\Delta s_i)+s_i\Delta x_i  & =0.
		\end{aligned}
	\]

	and since we know $s_i'=s_i+\Delta s_i$,
	\[
		\begin{aligned}
			x_is_i'+s_i\Delta x_i & =0                           \\
			x_is_i'               & =-s_i\Delta x_i              \\
			s_i'                  & =-\frac{s_i\Delta x_i}{x_i}.
		\end{aligned}
	\]

	since $x_i>0$, $s_i>0$, and $\Delta x_i>0$, we have
	\[
		s_i'<0
	\]

	\bigskip

	Now lets assume $\Delta s_i>0$. Prove $x_i'<0$

	Starting with the given equation,
	\[
		\begin{aligned}
			x_is_i+x_i\Delta s_i+s_i\Delta x_i & =0  \\
			s_i(x_i+\Delta x_i)+x_i\Delta s_i  & =0.
		\end{aligned}
	\]

	and since $x_i'=x_i+\Delta x_i$,
	\[
		\begin{aligned}
			s_ix_i'+x_i\Delta s_i & =0                           \\
			s_ix_i'               & =-x_i\Delta s_i              \\
			x_i'                  & =-\frac{x_i\Delta s_i}{s_i}.
		\end{aligned}
	\]

	knowing that  $x_i>0$, $s_i>0$, and $\Delta s_i>0$ we can state:
	\[
		x_i'<0
	\]

	If Newton's Method increases either $x_i$ or $s_i$,
	the full Newton step forces the other variable to become
	negative, moving outside the feasible region.

\end{homeworkProblem}

\begin{homeworkProblem}[A18]
	Adding slack variables $s_1,s_2$ to the dual and
	using $\mu=\frac13$, we can build the following system:
	\[
		\begin{aligned}
			x_1+2x_2 & =2       &  & (1)  \\
			y_1+s_1  & =3       &  & (2)  \\
			2y_1+s_2 & =4       &  & (3)  \\
			x_1s_1   & =\frac13 &  & (4)  \\
			x_2s_2   & =\frac13 &  & (5).
		\end{aligned}
	\]

	From (3) $-\,2(2)$, we obtain
	\[
		s_2-2s_1=-2,
	\]
	so that
	\[
		s_2=2s_1-2.
	\]

	Using (4) and (5), we can express $x_1$ and $x_2$
	in terms of $s_1$:
	\[
		x_1=\frac{1}{3s_1},
		\qquad
		x_2=\frac{1}{3s_2}
		=\frac{1}{6s_1-6}.
	\]

	Now substituting into (1),
	\[
		\begin{aligned}
			x_1+2x_2
			                                & =\frac{1}{3s_1}+\frac{2}{6s_1-6}=2 \\
			\frac{1}{3s_1}+\frac{1}{3s_1-3} & =2                                 \\
			\frac{2s_1-1}{3s_1(s_1-1)}      & =2                                 \\
			2s_1-1                          & =6s_1(s_1-1)                       \\
			2s_1-1                          & =6s_1^2-6s_1                       \\
			6s_1^2-8s_1+1                   & =0.
		\end{aligned}
	\]

	then we can apply the quadratic formula to get
	\[
		\begin{aligned}
			s_1
			 & =\frac{8\pm\sqrt{64-24}}{12} \\
			 & =\frac{8\pm\sqrt{40}}{12}    \\
			 & =\frac{4\pm\sqrt{10}}6.
		\end{aligned}
	\]

	But since $s_2=2s_1-2>0$, we require $s_1>1$ and we only take the positive root:
	\[
		s_1=\frac{4+\sqrt{10}}6
	\]

	Now we can use this found value in (4),
	\[
		\begin{aligned}
			x_1
			 & =\frac{1}{3s_1}               \\
			 & =\frac{2}{4+\sqrt{10}}        \\
			 & =\frac{2(4-\sqrt{10})}{16-10} \\
			 & =\frac{4-\sqrt{10}}3.
		\end{aligned}
	\]

	And using (1),
	\[
		\begin{aligned}
			x_2
			 & =\frac{2-x_1}{2}                 \\
			 & =\frac{2-\frac{4-\sqrt{10}}3}{2} \\
			 & =\frac{2+\sqrt{10}}6
		\end{aligned}
	\]

	Using (2),
	\[
		\begin{aligned}
			y_1
			 & =3-s_1                 \\
			 & =3-\frac{4+\sqrt{10}}6 \\
			 & =\frac{14-\sqrt{10}}6
		\end{aligned}
	\]

	And now we can solve \( s_{2} \),
	\[
		\begin{aligned}
			s_2
			 & =2s_1-2                              \\
			 & =2\left(\frac{4+\sqrt{10}}6\right)-2 \\
			 & =\frac{\sqrt{10}-2}3
		\end{aligned}
	\]

	Since all $x_i$ and $s_i$ are positive, and the five
	equations hold. We can state that the point on the central
	path with $\mu=\frac13$ is

	\[
		\begin{aligned}
			(x_1,x_2,y_1,s_1,s_2)
			=\bigg(
			 & \frac{4-\sqrt{10}}3,\,
			\frac{2+\sqrt{10}}6,
			 & \frac{14-\sqrt{10}}6,\,
			\frac{4+\sqrt{10}}6,\,
			\frac{\sqrt{10}-2}3
			\bigg).
		\end{aligned}
	\]
\end{homeworkProblem}
\end{document}
